\subsection{Code Evaluation}
\label{appendix:code-evaluation}

\subsubsection{[33] Orchestrator}

\paragraph{Modularity}

[33] demonstrated the use of modularity through the implementation of functions to separate different aspects of player behavior.
In the player script, horizontal movement, jumping, and gravity were handled in distinct functions, improving structural clarity.
Similarly, the instantiation of bullets in the shotgun script was separated into its own function.

The participant also made use of variables to parameterize behavior, such as controlling jump force and run speed, using clear and descriptive naming.
Within the bullet script, variables were used to represent speed and direction; however, these were given less descriptive names, reducing readability.
In the shotgun script, variables were used in accordance with the provided instructions.




\begin{itemize}
	\item \textbf{Score:}
	      \begin{tabular}{ll}
		      \textbf{Frequency:} & Often       \\
		      \textbf{Quality:}   & Appropriate \\
	      \end{tabular}
\end{itemize}

\paragraph{Instantiation}

[33] made use of instantiation for the creation of bullets, where instances were generated as part of the shooting mechanic.
The implementation was completed for this feature.
However, due to not finishing the workshop within the allotted time, instantiation was not applied to other relevant elements, such as enemy (ghost) spawning.


\begin{itemize}
	\item \textbf{Score:}
	      \begin{tabular}{ll}
		      \textbf{Frequency:} & Sometimes            \\
		      \textbf{Quality:}   & Somewhat Appropriate \\
	      \end{tabular}
\end{itemize}

\paragraph{State management}


[33] demonstrated the use of state management through the implementation of a pseudo state machine within the character controller, where input conditions were used to represent transitions between states.

A more explicit state machine was partially implemented for the shotgun mechanic; however, this remained incomplete, with only two states (idle and shooting) defined due to time constraints.

Changes to variable values were present and handled in a somewhat appropriate manner, although the structure did not consistently follow a clearly defined state management pattern.


\begin{itemize}
	\item \textbf{Score:}
	      \begin{tabular}{ll}
		      \textbf{Frequency:} & Sometimes            \\
		      \textbf{Quality:}   & Somewhat Appropriate \\
	      \end{tabular}
\end{itemize}

\paragraph{Events}


[33] made limited use of event handling, with a single signal implemented in the shotgun script.
This signal was triggered when the shooting animation reached its end, initiating a transition from the shooting state back to the idle state within the state machine.



\begin{itemize}
	\item \textbf{Score:}
	      \begin{tabular}{ll}
		      \textbf{Frequency:} & Sometimes   \\
		      \textbf{Quality:}   & Appropriate \\
	      \end{tabular}
\end{itemize}


\subsubsection{[11] Orchestrator}

\paragraph{Modularity}

[11] made use of modularity through both variables and functions across multiple scripts, but not always in the most optimal way.
In the [Player] script the tester created the function \texttt{jump\_move} which has the functionality for moving while in a jump, and the rest of the movement, gravity and jump logic is still in the main update loop.

Furthermore, variables such as \texttt{speed}, \texttt{jump\_force}, and \texttt{shootReady} are clearly used to structure behavior and store state and variable values.

\begin{itemize}
	\item \textbf{Score:}
	      \begin{tabular}{ll}
		      \textbf{Frequency:} & Sometimes     \\
		      \textbf{Quality:}   & Inappropriate \\
	      \end{tabular}
\end{itemize}

\paragraph{Instantiation}

[11] implemented instantiation in a clear and mostly correct way.
But didn't have time to finish the implementation.
In the shotgun script the bullets are properly instantiated and added to the scene whenever the input button is pressed, and it is a state which allows for shooting, and the bullet is then added to the scene.
For the ghost it instantiates the ghost each update but does not add it to the scene tree.

\begin{itemize}
	\item \textbf{Score:}
	      \begin{tabular}{ll}
		      \textbf{Frequency:} & Sometimes            \\
		      \textbf{Quality:}   & Somewhat Appropriate \\
	      \end{tabular}
\end{itemize}

\paragraph{State management}

[11] implemented two forms of state management: one for player movement and one for the shotgun mechanic.
However, the shotgun state machine was not fully implemented.

The player movement logic exhibited limitations in state transitions.
The player could only transition to the jumping state when not moving horizontally, and once in the jumping state, horizontal movement was no longer possible.

For the shotgun mechanic, only two states (shooting and idle) were defined.
The transition back to the idle state was handled using a delay, rather than an event signaling the completion of the shooting animation.

\begin{itemize}
	\item \textbf{Score:}
	      \begin{tabular}{ll}
		      \textbf{Frequency:} & Sometimes            \\
		      \textbf{Quality:}   & Somewhat Appropriate \\
	      \end{tabular}
\end{itemize}

\paragraph{Events}

[11] utilized event handling once to manage collisions between bullets and ghosts.
A signal was connected to the bullet script, which is triggered whenever the bullet’s collider detects a collision.
Upon activation, the corresponding function checks whether the collision involves a ghost entity.
If this condition is satisfied, both the ghost and the bullet are removed from the scene.


\begin{itemize}
	\item \textbf{Score:}
	      \begin{tabular}{ll}
		      \textbf{Frequency:} & Sometimes   \\
		      \textbf{Quality:}   & Appropriate \\
	      \end{tabular}
\end{itemize}


\subsubsection{22: Scratch}

\paragraph{Modularity}

[22] as mentioned before used custom blocks/functions right away, but the way they were placed in Scratch left some to be desired.
The Scratch code window can quickly become cluttered, and that happened to [22].
For example [22] made a ``movement right'' and a ``movement left'' function, but had the conditional key press statement outside the function.
This conditional "if" if added to the function could have reduced a lot of visual clutter.

Was quick to use it, and used it a lot in the ``knight'' part of the workshop.
The usage of functions did decline for other sprites, as they did not need as much functionality like the knight.

The functions created had a purpose and a reusability to them.
But as mentioned [22] did not catch that they could add the conditionals into most functions.
In the case of the movement functions for left and right, it would have been possible to combine into a singular function, and have the key pressed conditionals within it.
That would have been a great option to reduce clutter and would still fit in with existing logic [22] made.



\begin{itemize}
	\item \textbf{Score:}
	      \begin{tabular}{ll}
		      \textbf{Frequency:} & Usually     \\
		      \textbf{Quality:}   & Appropriate \\
	      \end{tabular}
\end{itemize}

\paragraph{Instantiation}

[22] used instantiation for both the ``ghost'' and ``bullet'' sprites as per the gamedev instructions.
Those two sprites were the only intended places to use instantiation. The final implementation of the ghost spawning mechanic was differnt from the provided intructions. [22] made a wave based spawning system, where X amount of ghosts would spawn per X wave.


\begin{itemize}
	\item \textbf{Score:}
	      \begin{tabular}{ll}
		      \textbf{Frequency:} & Often                \\
		      \textbf{Quality:}   & Somewhat Appropriate \\
	      \end{tabular}
\end{itemize}

\paragraph{State management}

[22] without instructions or help, used the built-in tools for Scratch to make a new costume of the knight, duplicating the original sprite and flipping it. Costume changes were added to left and right movement functions,  so the knight would be looking the way they were going. The left and right movement was two seperate functions and could be seen as states.



\begin{itemize}
	\item \textbf{Score:}
	      \begin{tabular}{ll}
		      \textbf{Frequency:} & Sometimes            \\
		      \textbf{Quality:}   & Somewhat Appropriate \\
	      \end{tabular}
\end{itemize}

\paragraph{Events}

[22] used basic event handling concepts for the keyboard inputs and collision detection. They used the built "touches" block to detect collisions, that was the right thing to do in Scrach.
There where attempts to use the "broadcast", but they were unsuccessful. Indicating either a lack of experience with event-driven as a concept or programming in genreal.
	[22] used for loops and if statements to check for key pressing events for Knight movement.
If [22] had used the event block for key presses, the movement would have been poor.


\begin{itemize}
	\item \textbf{Score:}
	      \begin{tabular}{ll}
		      \textbf{Frequency:} & Often                \\
		      \textbf{Quality:}   & Somewhat Appropriate \\
	      \end{tabular}
\end{itemize}



\renewcommand{\arraystretch}{1.1}

\begin{table}[h]
	\centering
	\footnotesize
	\begin{tabular}{
			>{\raggedright\arraybackslash}p{0.18\columnwidth}
			>{\raggedright\arraybackslash}p{0.20\columnwidth}
			>{\raggedright\arraybackslash}p{0.20\columnwidth}
			>{\raggedright\arraybackslash}p{0.20\columnwidth}
		}
		\hline
		\textbf{Concept} & \textbf{33}                                                         & \textbf{11} & \textbf{22} \\
		\hline

		Modularity
		                 & \textbf{Freq.:} Often \newline \textbf{Qual.:} Appropriate
		                 & \textbf{Freq.:} Often \newline \textbf{Qual.:} Somewhat
		                 & \textbf{Freq.:} Usually \newline \textbf{Qual.:} Appropriate                                    \\

		Instantiation
		                 & \textbf{Freq.:} Sometimes \newline \textbf{Qual.:} Somewhat Approp.
		                 & \textbf{Freq.:} Sometimes \newline \textbf{Qual.:} Somewhat Approp.
		                 & \textbf{Freq.:} Often \newline \textbf{Qual.:} Somewhat Approp.                                 \\

		State management
		                 & \textbf{Freq.:} Often \newline \textbf{Qual.:} Somewhat Approp.
		                 & \textbf{Freq.:} Often \newline \textbf{Qual.:} Somewhat Approp.
		                 & \textbf{Freq.:} Sometimes \newline \textbf{Qual.:} Somewhat Approp.                             \\

		Events
		                 & \textbf{Freq.:} Sometimes \newline \textbf{Qual.:} Appropriate
		                 & \textbf{Freq.:} Sometimes \newline \textbf{Qual.:} Somewhat Approp.
		                 & \textbf{Freq.:} Often \newline \textbf{Qual.:} Somewhat Approp.                                 \\

		\hline
	\end{tabular}
	\caption{Summary of concept evaluation}\label{tab:ArtifactSummarytable}
\end{table}